\documentclass[12pt, a4paper]{article}

% ── Packages ──────────────────────────────────────────────────────────────────
\usepackage[T1]{fontenc}
\usepackage[utf8]{inputenc}
\usepackage[margin=2.4cm, top=2.8cm, bottom=2.8cm]{geometry}
\usepackage{lmodern}
\usepackage{microtype}
\usepackage{parskip}
\usepackage{titlesec}
\usepackage{xcolor}
\usepackage{mdframed}
\usepackage{enumitem}
\usepackage{booktabs}
\usepackage{array}
\usepackage{fancyhdr}
\usepackage{hyperref}
\usepackage{tikz}
\usepackage{pgfplots}
\pgfplotsset{compat=1.18}
\usetikzlibrary{shapes.geometric, shapes.rounded, arrows.meta, positioning,
                fit, backgrounds, shadows.blur, calc}

% ── Colors ────────────────────────────────────────────────────────────────────
\definecolor{accent}{RGB}{94, 106, 210}
\definecolor{accentlight}{RGB}{220, 223, 245}
\definecolor{lightgray}{RGB}{245, 246, 248}
\definecolor{darkgray}{RGB}{50, 55, 70}
\definecolor{mutedtext}{RGB}{110, 115, 135}
\definecolor{csource}{RGB}{167, 139, 250}   % purple  - sources
\definecolor{cengine}{RGB}{94, 106, 210}    % indigo  - engine
\definecolor{coutput}{RGB}{52, 211, 153}    % emerald - output
\definecolor{cdanger}{RGB}{239, 68, 68}     % red     - problem
\definecolor{cokay}{RGB}{34, 197, 94}       % green   - solution
\definecolor{camber}{RGB}{245, 158, 11}     % amber   - warning / phase

% ── Hyperref ──────────────────────────────────────────────────────────────────
\hypersetup{colorlinks=true, linkcolor=accent, urlcolor=accent,
            pdftitle={dedomena Product Overview}}

% ── Section styles ────────────────────────────────────────────────────────────
\titleformat{\section}
  {\large\bfseries\color{darkgray}}{}{}{}
  [\vspace{-0.3em}\textcolor{accent}{\rule{\linewidth}{1.5pt}}\vspace{0.4em}]
\titleformat{\subsection}{\normalsize\bfseries\color{darkgray}}{}{}{}
\titlespacing{\section}{0pt}{1.8em}{0.6em}
\titlespacing{\subsection}{0pt}{1.2em}{0.3em}

% ── Header / footer ───────────────────────────────────────────────────────────
\pagestyle{fancy}
\fancyhf{}
\fancyhead[L]{\small\textcolor{mutedtext}{dedomena \textsigma}}
\fancyhead[R]{\small\textcolor{mutedtext}{Product Overview}}
\fancyfoot[C]{\small\textcolor{mutedtext}{\thepage}}
\renewcommand{\headrulewidth}{0pt}

% ── Callout boxes ─────────────────────────────────────────────────────────────
\newmdenv[backgroundcolor=lightgray, linecolor=accent, linewidth=3pt,
          topline=false, bottomline=false, rightline=false,
          innerleftmargin=14pt, innerrightmargin=14pt,
          innertopmargin=10pt, innerbottommargin=10pt,
          skipabove=10pt, skipbelow=10pt]{callout}

\newmdenv[backgroundcolor=accent!7, linecolor=accent!35, linewidth=1pt,
          roundcorner=5pt, innerleftmargin=14pt, innerrightmargin=14pt,
          innertopmargin=10pt, innerbottommargin=10pt,
          skipabove=10pt, skipbelow=10pt]{highlight}

% ── TikZ styles ───────────────────────────────────────────────────────────────
\tikzset{
  % Generic rounded box
  box/.style={
    rectangle, rounded corners=6pt,
    draw=#1!60, fill=#1!10, text=#1!80!black,
    minimum height=1.0cm, text centered,
    font=\small\bfseries, inner sep=8pt,
    drop shadow={shadow xshift=1pt, shadow yshift=-1pt,
                 fill=#1!20, opacity=0.5}
  },
  % Arrow
  arr/.style={-{Stealth[length=6pt]}, thick, color=#1, shorten >=2pt, shorten <=2pt},
  % Small pill label on arrows
  alabel/.style={font=\scriptsize\itshape, text=mutedtext,
                 fill=white, inner sep=2pt, midway, sloped=false},
}

% ── Document ──────────────────────────────────────────────────────────────────
\begin{document}

% ── Cover block ───────────────────────────────────────────────────────────────
\thispagestyle{empty}
\begin{center}
  \vspace*{1.4cm}
  {\fontsize{36}{42}\selectfont\bfseries\color{darkgray} dedomena}
  \hspace{0.4em}
  {\fontsize{36}{42}\selectfont\color{accent!45}\textsigma}

  \vspace{0.6em}
  {\large\color{mutedtext} A smarter way to work with your data}

  \vspace{1.6em}
  \textcolor{accent!30}{\rule{7cm}{0.6pt}}

  \vspace{1em}
  {\small\color{mutedtext}
    Product Overview \quad\textbullet\quad Internal Document \quad\textbullet\quad \today
  }
  \vspace{0.6cm}

  % ── Cover diagram: the one-line pitch ──────────────────────────────────────
  \begin{tikzpicture}[node distance=0.8cm and 1.1cm]
    \node[box=csource, minimum width=3.2cm] (src) {Your Data};
    \node[box=cengine, minimum width=3.2cm, right=of src] (eng) {dedomena};
    \node[box=coutput, minimum width=3.2cm, right=of eng] (out) {Clear Answer};

    \draw[arr=csource] (src) -- (eng)
      node[alabel, above] {connects};
    \draw[arr=cengine] (eng) -- (out)
      node[alabel, above] {delivers};

    \node[font=\scriptsize\color{mutedtext}, below=0.35cm of src]
      {100+ sources};
    \node[font=\scriptsize\color{mutedtext}, below=0.35cm of eng]
      {AI reasoning};
    \node[font=\scriptsize\color{mutedtext}, below=0.35cm of out]
      {in seconds};
  \end{tikzpicture}

  \vspace{1.6cm}
\end{center}
\newpage

% ── 1. What is dedomena ───────────────────────────────────────────────────────
\section{What Is dedomena?}

dedomena is a web platform that lets anyone in the organisation ask questions
about business data in plain English and get clear, sourced answers in seconds.

No SQL. No Excel formulas. No waiting for an analyst to run a report.
You connect your data sources once, type your question, and the platform does the rest.

\begin{callout}
  \textbf{Think of it like having a very knowledgeable assistant who has read every
  spreadsheet, every database table, and every report in the company, and who is
  always ready to answer a question instantly.}
\end{callout}


% ── 2. The problem ────────────────────────────────────────────────────────────
\section{The Problem We Are Solving}

\subsection{The old way is too slow}

Getting an answer from data today involves too many steps and too many people.
The diagram below shows what a typical request looks like right now.

\vspace{0.6em}
\begin{center}
\begin{tikzpicture}[node distance=0.5cm and 0.7cm,
    step/.style={rectangle, rounded corners=4pt, draw=cdanger!50,
                 fill=cdanger!8, text=cdanger!80!black, font=\small,
                 minimum height=0.85cm, minimum width=3.0cm,
                 text centered, inner sep=6pt},
    wait/.style={font=\scriptsize\itshape\color{mutedtext},
                 fill=white, inner sep=2pt}]

  \node[step] (s1) {Manager needs a number};
  \node[step, right=of s1] (s2) {Emails analyst};
  \node[step, right=of s2] (s3) {Analyst finds data};
  \node[step, below=of s3] (s4) {Analyst writes query};
  \node[step, left=of s4]  (s5) {Formats a report};
  \node[step, left=of s5]  (s6) {Manager gets answer};

  \draw[arr=cdanger] (s1)--(s2) node[wait, above]{};
  \draw[arr=cdanger] (s2)--(s3) node[wait, above]{1--2 days};
  \draw[arr=cdanger] (s3)--(s4);
  \draw[arr=cdanger] (s4)--(s5) node[wait, above]{hours};
  \draw[arr=cdanger] (s5)--(s6);

  \node[font=\scriptsize\bfseries\color{cdanger},
        below=0.7cm of s5, xshift=1.5cm]
    {Total time: 1 to 3 days for a single question};
\end{tikzpicture}
\end{center}
\vspace{0.4em}

This is slow. It creates a bottleneck on a small number of technical people,
and it means decisions are made without information that already exists.

\subsection{Data lives in too many places}

Beyond speed, a second problem is that data is scattered. No single tool gives a
combined view, which means every cross-system question requires manual work.

\vspace{0.6em}
\begin{center}
\begin{tikzpicture}[
    silo/.style={rectangle, rounded corners=5pt,
                 draw=mutedtext!50, fill=lightgray,
                 font=\small, minimum width=2.2cm, minimum height=0.9cm,
                 text centered, inner sep=6pt, text=darkgray},
    question/.style={ellipse, draw=cdanger!60, fill=cdanger!8,
                     font=\small\itshape, text=cdanger!80!black,
                     inner sep=8pt}]

  \node[silo] (crm)  at (0,0)    {CRM};
  \node[silo] (fin)  at (2.6,0)  {Finance};
  \node[silo] (ops)  at (5.2,0)  {Operations};
  \node[silo] (hr)   at (7.8,0)  {HR};
  \node[silo] (files)at (3.9,-1.6){Spreadsheets};

  \node[question] (q) at (3.9, -3.2)
    {"Which region grew fastest this quarter?"};

  \foreach \n in {crm, fin, ops, hr, files}{
    \draw[thick, dashed, color=mutedtext!50] (\n) -- (q);
  }

  \node[font=\scriptsize\color{mutedtext}, below=0.25cm of q]
    {No single tool can answer this. Someone has to pull from all of them manually.};
\end{tikzpicture}
\end{center}
\vspace{0.4em}


% ── 3. How we will use dedomena ───────────────────────────────────────────────
\section{How We Will Use dedomena}

\subsection{The new way: one question, one answer}

With dedomena, the same request that used to take days now takes seconds.
The entire workflow collapses into three steps.

\vspace{0.6em}
\begin{center}
\begin{tikzpicture}[node distance=1.0cm and 1.4cm,
    step/.style={rectangle, rounded corners=6pt,
                 draw=cokay!60, fill=cokay!8, text=cokay!80!black,
                 font=\small\bfseries, minimum height=1.0cm,
                 minimum width=3.6cm, text centered, inner sep=8pt,
                 drop shadow={shadow xshift=1pt, shadow yshift=-1pt,
                              fill=cokay!20, opacity=0.4}}]

  \node[step] (a) {Ask a question};
  \node[step, right=of a] (b) {dedomena reads all sources};
  \node[step, right=of b] (c) {Answer delivered};

  \draw[arr=cokay] (a)--(b)
    node[midway, above, font=\scriptsize\itshape\color{mutedtext}]{instantly};
  \draw[arr=cokay] (b)--(c)
    node[midway, above, font=\scriptsize\itshape\color{mutedtext}]{seconds};

  \node[font=\scriptsize\bfseries\color{cokay},
        below=0.7cm of b]{Total time: under 30 seconds};
\end{tikzpicture}
\end{center}
\vspace{0.4em}

\subsection{Step one: connect our data sources}

We connect the tools and databases we already use. This is a one-time setup
that takes a few minutes per source. Once connected, dedomena keeps the data
ready for questions. Nothing is sent to a third party. Everything stays in our
environment.

\vspace{0.8em}
\begin{center}
\begin{tikzpicture}[
    cat/.style={rectangle, rounded corners=5pt,
                draw=csource!50, fill=csource!8,
                font=\small\bfseries, text=csource!80!black,
                minimum width=3.0cm, minimum height=0.85cm,
                text centered, inner sep=8pt},
    eng/.style={rectangle, rounded corners=8pt,
                draw=cengine!60, fill=cengine!12,
                font=\normalsize\bfseries, text=cengine!80!black,
                minimum width=2.6cm, minimum height=3.8cm,
                text centered, inner sep=10pt,
                drop shadow={shadow xshift=1.5pt, shadow yshift=-1.5pt,
                             fill=cengine!25, opacity=0.5}},
    arr/.style={-{Stealth[length=5pt]}, thick, color=csource!60}]

  % Source nodes (left column)
  \node[cat] (crm)  at (0, 2.4)  {CRM};
  \node[cat] (db)   at (0, 1.2)  {Databases};
  \node[cat] (files)at (0, 0.0)  {Files};
  \node[cat] (cloud)at (0,-1.2)  {Cloud Storage};
  \node[cat] (fin)  at (0,-2.4)  {Finance};

  % dedomena engine (center)
  \node[eng] (eng)  at (4.2, 0)  {dedomena};

  % Answer node (right)
  \node[rectangle, rounded corners=6pt,
        draw=coutput!60, fill=coutput!10,
        font=\small\bfseries, text=coutput!80!black,
        minimum width=2.8cm, minimum height=1.2cm,
        text centered, inner sep=10pt,
        drop shadow={shadow xshift=1pt, shadow yshift=-1pt,
                     fill=coutput!25, opacity=0.4}]
       (out) at (8.0, 0) {Your Answer};

  % Arrows from sources to engine
  \foreach \s in {crm, db, files, cloud, fin}{
    \draw[arr] (\s) -- (eng);
  }

  % Arrow from engine to answer
  \draw[-{Stealth[length=6pt]}, thick, color=cengine!70] (eng) -- (out)
    node[midway, above, font=\scriptsize\itshape\color{mutedtext}]{seconds};

  % Labels
  \node[font=\scriptsize\color{mutedtext}, below=0.3cm of fin]{100+ sources supported};
\end{tikzpicture}
\end{center}
\vspace{0.4em}

\subsection{Step two: ask questions in plain English}

Any team member with access to the platform opens dedomena, types a question,
and receives an answer. Here are realistic examples of what we would actually type:

\begin{highlight}
  \textit{"What were our top five accounts by revenue last quarter?"}

  \vspace{0.4em}
  \textit{"Show me a summary of all open support tickets created this month."}

  \vspace{0.4em}
  \textit{"Which product line has the highest return rate according to our sales data?"}

  \vspace{0.4em}
  \textit{"Compare our headcount from January to today and flag any department that
           grew by more than 20 percent."}
\end{highlight}

dedomena reads across all connected sources at once, finds the relevant information,
and writes a clear answer with references to where the data came from.

\subsection{Step three: act on the results}

The platform does not just return a number. It explains what it found, flags
anything unusual, and presents the information as a summary, a table, or a list,
depending on what is most useful. The person asking has everything they need to
make a decision without opening a single spreadsheet.


% ── 4. Who benefits ───────────────────────────────────────────────────────────
\section{Who Benefits and How}

\vspace{0.4em}
\begin{center}
\begin{tikzpicture}[
    role/.style={rectangle, rounded corners=6pt,
                 draw=cengine!40, fill=cengine!7,
                 font=\small\bfseries, text=darkgray,
                 minimum width=2.6cm, minimum height=0.85cm,
                 text centered, inner sep=8pt},
    gain/.style={rectangle, rounded corners=4pt,
                 draw=coutput!40, fill=coutput!7,
                 font=\scriptsize, text=darkgray,
                 minimum width=5.2cm, text width=5.0cm,
                 inner sep=7pt},
    conn/.style={-{Stealth[length=5pt]}, thick, color=mutedtext!60}]

  % Roles (left)
  \node[role] (mgr)  at (0,  3.0) {Management};
  \node[role] (sales)at (0,  1.5) {Sales};
  \node[role] (ops)  at (0,  0.0) {Operations};
  \node[role] (fin)  at (0, -1.5) {Finance};
  \node[role] (data) at (0, -3.0) {Analysts};

  % Gains (right)
  \node[gain] (gmgr)  at (7.4,  3.0)
    {Instant answers to performance questions. No more waiting for reports.};
  \node[gain] (gsales)at (7.4,  1.5)
    {Customer history and account health at their fingertips.};
  \node[gain] (gops)  at (7.4,  0.0)
    {A combined view of inventory, suppliers, and process data across systems.};
  \node[gain] (gfin)  at (7.4, -1.5)
    {Cross-reference financials and flag budget issues without formulas.};
  \node[gain] (gdata) at (7.4, -3.0)
    {Free from routine data requests. Focus on deeper analysis.};

  % Connecting arrows
  \foreach \r/\g in {mgr/gmgr, sales/gsales, ops/gops, fin/gfin, data/gdata}{
    \draw[conn] (\r) -- (\g);
  }
\end{tikzpicture}
\end{center}
\vspace{0.4em}


% ── 5. What makes it different ────────────────────────────────────────────────
\section{What Makes It Different}

\vspace{0.6em}
\begin{center}
\begin{tikzpicture}[
    col/.style={rectangle, rounded corners=6pt,
                draw=#1!50, fill=#1!8,
                minimum width=5.6cm, text width=5.4cm,
                inner sep=12pt, font=\small, text=#1!80!black},
    head/.style={font=\small\bfseries, text=#1!80!black}]

  % Column headers
  \node[font=\small\bfseries\color{cdanger!80}] at (0, 0) {Without dedomena};
  \node[font=\small\bfseries\color{cokay!80}]   at (7.4, 0) {With dedomena};

  % Divider
  \draw[thick, dashed, color=mutedtext!40] (3.7, 0.5) -- (3.7, -7.8);

  % Comparison rows
  \foreach \y/\bad/\good in {
    -1.0/{Waiting 1 to 3 days for an answer}/{Answer in under 30 seconds},
    -2.6/{One dashboard at a time}/{Reads all sources at once},
    -4.2/{Need to know SQL or Excel}/{Plain English questions only},
    -5.8/{Analyst is the bottleneck}/{Any team member can ask},
    -7.4/{Data spread across 10 tools}/{One place for all questions}
  }{
    \node[col=cdanger, align=left] at (0, \y) {\bad};
    \node[col=cokay,   align=left] at (7.4, \y) {\good};
    \node[font=\large\color{accent}] at (3.7, \y) {$\rightarrow$};
  }

\end{tikzpicture}
\end{center}
\vspace{0.4em}

\subsection{The data stays with us}

This is critical from a compliance standpoint. dedomena does not upload your raw
data to any external server. When a question is asked, only the portions of data
that are relevant to that specific question are sent to the AI to generate
the answer. The rest stays in place, inside our environment.


% ── 6. What we get as an organisation ────────────────────────────────────────
\section{What We Get as an Organisation}

\begin{enumerate}[leftmargin=1.6em, itemsep=8pt]
  \item \textbf{Faster decision-making.}
  Questions that used to take days take seconds.
  Teams act on information while it is still relevant.

  \item \textbf{Less dependency on technical staff for routine questions.}
  Analysts are freed from handling basic data requests.
  They spend their time on work that actually requires their skills.

  \item \textbf{One place to ask questions about all of our data.}
  Instead of logging into five different tools and exporting five spreadsheets,
  everyone queries through one interface.

  \item \textbf{A full record of what was asked and when.}
  Every question is logged with the sources it used, creating a natural
  audit trail for any data-driven decision.

  \item \textbf{Better data engagement across the organisation.}
  When people can ask questions directly, they engage with data more often.
  Over time, this builds a culture where decisions are grounded in evidence.
\end{enumerate}


% ── 7. Rollout plan ───────────────────────────────────────────────────────────
\section{Suggested Rollout Plan}

\vspace{0.8em}
\begin{center}
\begin{tikzpicture}[
    phase/.style={rectangle, rounded corners=6pt,
                  draw=#1!60, fill=#1!10,
                  minimum width=4.0cm, text width=3.8cm,
                  inner sep=12pt, font=\small, text=darkgray,
                  drop shadow={shadow xshift=1pt, shadow yshift=-1pt,
                               fill=#1!20, opacity=0.45}},
    badge/.style={circle, draw=#1!60, fill=#1!15,
                  font=\small\bfseries, text=#1!80!black,
                  minimum size=0.9cm},
    phaseconn/.style={-{Stealth[length=6pt]}, thick, color=mutedtext!50}]

  % Phase nodes
  \node[phase=camber] (p1) at (0, 0) {
    \textbf{Phase 1: Pilot}\\[4pt]
    Connect 2 or 3 key sources.\\
    Run a 2-week trial with managers and team leads.\\
    Collect feedback.
  };

  \node[phase=cengine] (p2) at (5.2, 0) {
    \textbf{Phase 2: Rollout}\\[4pt]
    Connect all data sources.\\
    Open access to all teams.\\
    Run a short onboarding session.
  };

  \node[phase=coutput] (p3) at (10.4, 0) {
    \textbf{Phase 3: Optimise}\\[4pt]
    Review the most common questions.\\
    Build standing summaries for recurring needs.
  };

  % Timeline line
  \draw[thick, color=mutedtext!40] (-2, -1.8) -- (12.4, -1.8);

  % Phase markers on timeline
  \node[badge=camber] at (0, -1.8)  {1};
  \node[badge=cengine] at (5.2, -1.8) {2};
  \node[badge=coutput] at (10.4, -1.8) {3};

  % Time labels
  \node[font=\scriptsize\color{mutedtext}] at (0, -2.55)  {Weeks 1 to 2};
  \node[font=\scriptsize\color{mutedtext}] at (5.2, -2.55) {Month 2};
  \node[font=\scriptsize\color{mutedtext}] at (10.4, -2.55){Month 3 onwards};

  % Connectors between phases
  \draw[phaseconn] (p1) -- (p2);
  \draw[phaseconn] (p2) -- (p3);

\end{tikzpicture}
\end{center}
\vspace{0.4em}


% ── 8. Summary ────────────────────────────────────────────────────────────────
\section{Summary}

dedomena gives our organisation the ability to ask any question about our data and
get a clear, reliable answer in seconds, without writing code, without waiting for
a report, and without needing to know where the data lives.

We will use it to reduce the time between a question and a decision, to remove the
bottleneck that currently sits on our analytics team, and to give every member of
the team direct access to the information they need to do their job well.

\begin{callout}
  \textbf{The goal is simple:} anyone in the organisation should be able to open
  dedomena, type a question about our business, and have a clear answer in front
  of them before their coffee gets cold.
\end{callout}

\vspace{2em}
\textcolor{mutedtext}{\small This document is intended for internal review. For a
live demonstration or to discuss integration with our existing tools, please reach
out to the product team.}

\end{document}
